\fiche{Solutions des dossiers}{Méthodes et conclusions}\par\addvspace{6mm}\Needspace{9\baselineskip}\souscours{1. Une serre à construire}\phantomsection\label{projet-cor-1}\pdfbookmark[1]{Une serre à construire}{bm-projet-cor-1}
\repere{projet-1}{Énoncé}
\begin{questions}
\item La longueur vaut $40-2x$ ; $A(x)=40x-2x^2$, avec $0<x<20$.
\item $A\,\!^{\prime}=40-4x$ : le maximum est en $x=10$, avec une longueur de $20$ m et une aire de $200$ m$^2$.
\item $40x-2x^2\geq150$ équivaut à $x^2-20x+75\leq0$, donc $5\leq x\leq15$.
\item Le schéma doit montrer le mur, deux largeurs $x$ et le côté $40-2x$. La solution respecte les $40$ m disponibles ; le modèle néglige l’épaisseur des matériaux.
\end{questions}
\par\addvspace{6mm}\Needspace{9\baselineskip}\souscours{2. Deux modèles de refroidissement}\phantomsection\label{projet-cor-2}\pdfbookmark[1]{Deux modèles de refroidissement}{bm-projet-cor-2}
\repere{projet-2}{Énoncé}
\begin{questions}
\item $T(t)=20+60\mathrm e^{-0{,}1t}$.
\item $u_{n+1}=0{,}9u_n+2$, $u_0=80$ ; donc $u_n=20+60(0{,}9)^n$.
\item $T(10)\simeq42{,}073$ et $u_{10}\simeq40{,}921$ ; Euler est ici plus bas d’environ $1{,}152$.
\item Réduire le pas d’Euler et comparer à la solution exacte. Une précision numérique meilleure ne corrige toutefois pas une hypothèse physique inadaptée.
\end{questions}
\par\addvspace{6mm}\Needspace{9\baselineskip}\souscours{3. Du produit à la somme}\phantomsection\label{projet-cor-3}\pdfbookmark[1]{Du produit à la somme}{bm-projet-cor-3}
\repere{projet-3}{Énoncé}
\begin{questions}
\item La propriété du logarithme népérien, divisée par $\ln10$, donne l’identité.
\item $\ell(6)\simeq0{,}77815$ et $\ell(12)=2\ell(2)+\ell(3)\simeq1{,}07918$.
\item Une primitive de $1/t$ est $\ln t$ : $\int_1^x\mathrm dt/t=\ln x$.
\item Les puissances d’une base, en progression géométrique, ont des logarithmes en progression arithmétique. Une autre construction du logarithme utilise l’aire sous l’hyperbole. Les deux approches conduisent à la même fonction après normalisation.
\end{questions}
\par\addvspace{6mm}\Needspace{9\baselineskip}\souscours{4. Une aire sans primitive usuelle}\phantomsection\label{projet-cor-4}\pdfbookmark[1]{Une aire sans primitive usuelle}{bm-projet-cor-4}
\repere{projet-4}{Énoncé}
\begin{questions}
\item Sur $[0,1]$, $\mathrm e^{-1}\leq\mathrm e^{-x^2}\leq1$, donc $\mathrm e^{-1}\leq I\leq1$.
\item La fonction décroît : $D_n=\frac1n\sum_{k=1}^n\mathrm e^{-(k/n)^2}\leq I\leq G_n=\frac1n\sum_{k=0}^{n-1}\mathrm e^{-(k/n)^2}$. La largeur vaut $(1-\mathrm e^{-1})/n$.
\item Tirer $N$ couples indépendants $(U,V)$ uniformes dans le carré ; compter les $V\leq\mathrm e^{-U^2}$ et rendre ce nombre divisé par $N$.
\item Les rectangles donnent un encadrement déterministe sous les hypothèses de monotonie. Monte-Carlo donne une estimation aléatoire dont la fluctuation doit être étudiée ; il ne donne pas à lui seul une borne certaine.
\end{questions}
\par\addvspace{6mm}\Needspace{9\baselineskip}\souscours{5. Mesurer une inégalité}\phantomsection\label{projet-cor-5}\pdfbookmark[1]{Mesurer une inégalité}{bm-projet-cor-5}
\repere{projet-5}{Énoncé}
\begin{questions}
\item Les deux vont de $0$ à $1$, sont croissantes, convexes et sous la diagonale : pour $B$, $L_B\,\!^{\prime}=(1+3x^2)/2>0$ et $L_B\,\!^{\prime\prime}=3x\geq0$.
\item $G_A=1/3$ ; $\int_0^1L_B=\tfrac12(1/2+1/4)=3/8$, donc $G_B=1/4$.
\item $L_B-L_A=x(x-1)^2/2\geq0$ : $B$ est plus proche de l’égalité au sens de Lorenz.
\item Non : les parts relatives restent identiques. Le niveau des revenus et leur répartition relative sont des informations distinctes.
\end{questions}
\par\addvspace{6mm}\Needspace{9\baselineskip}\souscours{6. Une information qui change la décision}\phantomsection\label{projet-cor-6}\pdfbookmark[1]{Une information qui change la décision}{bm-projet-cor-6}
\repere{projet-6}{Énoncé}
\begin{questions}
\item Les chemins défectueux ont probabilités $0{,}8(0{,}01)=0{,}008$ et $0{,}2(0{,}06)=0{,}012$ ; total $0{,}02$.
\item $0{,}012/0{,}02=0{,}6$.
\item $0{,}06p>0{,}01(1-p)$ équivaut à $p>1/7$.
\item Avant l’observation, $P(B)=0{,}2$ ; après observation d’un défaut, $P_D(B)=0{,}6$. Le taux plus élevé de défaut chez $B$ explique cette révision.
\end{questions}
\par\addvspace{6mm}\Needspace{9\baselineskip}\souscours{7. Contrôler une surréservation}\phantomsection\label{projet-cor-7}\pdfbookmark[1]{Contrôler une surréservation}{bm-projet-cor-7}
\repere{projet-7}{Énoncé}
\begin{questions}
\item $X\sim\mathcal B(12,0{,}8)$.
\item $P(X=11)+P(X=12)=12(0{,}8)^{11}(0{,}2)+(0{,}8)^{12}\simeq0{,}274878$.
\item Non : le risque vaut environ $27{,}49\,\%$, au-dessus de $5\,\%$. Une espérance de $9{,}6$ présents inférieure à $10$ ne suffit pas.
\item Des clients voyageant en groupe créent une dépendance ; des probabilités de présence différentes selon les personnes empêchent l’usage direct de cette binomiale.
\end{questions}
\par\addvspace{6mm}\Needspace{9\baselineskip}\souscours{8. Attendre sans vieillir ?}\phantomsection\label{projet-cor-8}\pdfbookmark[1]{Attendre sans vieillir ?}{bm-projet-cor-8}
\repere{projet-8}{Énoncé}
\begin{questions}
\item $E(T)=E(U)=50$ heures ; pour l’exponentielle $\lambda=0{,}02$.
\item $\mathrm e^{-0{,}02\times30}=\mathrm e^{-0{,}6}\simeq0{,}548812$.
\item $(60/100)/(90/100)=2/3$.
\item $P(T>30)=\mathrm e^{-0{,}6}$ : égalité sans mémoire. $P(U>30)=0{,}7\ne2/3$ : l’uniforme n’a pas cette propriété. Une même moyenne n’implique pas une même loi.
\end{questions}
\par\addvspace{6mm}\Needspace{9\baselineskip}\souscours{9. Une croissance observée}\phantomsection\label{projet-cor-9}\pdfbookmark[1]{Une croissance observée}{bm-projet-cor-9}
\repere{projet-9}{Énoncé}
\begin{questions}
\item Les rapports successifs valent $2$ : $y=2\times2^t$ est un modèle compatible.
\item $\ln y=\ln2+t\ln2$ ; l’ajustement logarithmique est exactement affine.
\item $y(4)=32$ et $y(20)=2\,097\,152$. Ce sont deux extrapolations ; la seconde est beaucoup plus éloignée des données et peut être dénuée de pertinence.
\item Non. Le temps indexe les observations ; le mécanisme doit être expliqué par le contexte. Quatre couples ne démontrent pas une loi universelle ni un lien causal.
\end{questions}
\fiche{Solutions des évaluations}{Méthodes et conclusions}\par\addvspace{6mm}\Needspace{9\baselineskip}\souscours{1. Évaluation 1 · Analyse}\phantomsection\label{evaluation-cor-1}\pdfbookmark[1]{Évaluation 1 · Analyse}{bm-evaluation-cor-1}
\repere{evaluation-1}{Énoncé}
\begin{questions}
\item $f(0)=100$ ; $f\,\!^{\prime}=-20\mathrm e^{-t/4}<0$ ; limite $20$. Barème : 1, 2 et 1 point.
\item $\mathrm e^{-t/4}\leq1/4$ donne $t\geq4\ln4\simeq5{,}545$. Domaine et transformation : 2 points ; résolution et conclusion : 2 points.
\item $y\,\!^{\prime}=-y/4+5$, $y(0)=100$. Vérifier en substituant ; 2 points pour chaque élément.
\item Une primitive est $20t-320\mathrm e^{-t/4}$ ; moyenne $20+80(1-\mathrm e^{-1})\simeq70{,}570$. Primitive : 2 points ; calcul et interprétation : 2 points.
\item $u_{n+1}=0{,}75u_n+5$, $u_0=100$ ; $u_1=80,u_2=65$. L’exact donne $f(1)\simeq82{,}304$, $f(2)\simeq68{,}522$. Euler sous-estime ici. Récurrence : 2 points ; valeurs et comparaison : 2 points.
\end{questions}
\par\addvspace{6mm}\Needspace{9\baselineskip}\souscours{2. Évaluation 2 · Probabilités}\phantomsection\label{evaluation-cor-2}\pdfbookmark[1]{Évaluation 2 · Probabilités}{bm-evaluation-cor-2}
\repere{evaluation-2}{Énoncé}
\begin{questions}
\item $d=0{,}7(0{,}02)+0{,}3(0{,}05)=0{,}029$. Arbre ou formule : 2 points ; valeur : 2 points.
\item $P_D(B)=0{,}015/0{,}029=15/29\simeq0{,}517241$. Bayes : 2 points ; calcul : 2 points.
\item $X\sim\mathcal B(10,0{,}029)$ ; $P(X\geq1)=1-0{,}971^{10}\simeq0{,}254938$. Loi et hypothèses : 2 points ; complément : 2 points.
\item Géométrique sur $\N^*$ de paramètre $0{,}029$ ; $E=1/0{,}029\simeq34{,}483$ ; $P(X>20)=0{,}971^{20}\simeq0{,}555117$. Loi : 1 point ; espérance : 1 point ; survie : 2 points.
\item $\lambda=1/30$ ; les deux probabilités valent $\mathrm e^{-1/2}\simeq0{,}606531$. Paramètre et première probabilité : 2 points ; absence de mémoire : 2 points.
\end{questions}
\par\addvspace{6mm}\Needspace{9\baselineskip}\souscours{3. Évaluation 3 · Statistique et inégalités}\phantomsection\label{evaluation-cor-3}\pdfbookmark[1]{Évaluation 3 · Statistique et inégalités}{bm-evaluation-cor-3}
\repere{evaluation-3}{Énoncé}
\begin{questions}
\item $G(2{,}5,4)$, $V_x=1{,}25$, $V_y=2{,}5$, covariance $1{,}75$. Un point par quantité ou groupe de moyennes.
\item $y=1{,}4x+0{,}5$ ; prévisions $5{,}4$ et $14{,}5$. La première interpole, la seconde extrapole au-delà de $[1,4]$. Coefficients : 2 points ; valeurs et critique : 2 points.
\item $r=1{,}75/\sqrt{3{,}125}\simeq0{,}989949$ : forte association affine sur ces données, sans preuve causale. Calcul : 2 points ; commentaire : 2 points.
\item $L(0)=0,L(1)=1$, $L\,\!^{\prime}=2x\geq0$, $L\,\!^{\prime\prime}=2>0$ et $x^2\leq x$ sur $[0,1]$. $L(0{,}6)=0{,}36$ : les $60\,\%$ les moins riches possèdent $36\,\%$ de la richesse. Vérifications : 3 points ; lecture : 1 point.
\item $G=1-2\int_0^1x^2\,\mathrm dx=1/3$. Une multiplication commune conserve les parts relatives, donc Lorenz et Gini. Calcul : 2 points ; explication : 2 points.
\end{questions}
